% SPDX-License-Identifier: Apache-2.0
\section{\code{txhdl\_macros}}
\label{sec:r-macros}

The procedural macros, written against \code{proc\_macro} alone so the
crate has no dependencies. \code{\#[derive(Transaction)]} and
\code{\#[derive(Bus)]} emit one empty \code{impl} each.
\code{\#[derive(Value)]} gives a struct of values, or a fieldless
enum, its width and its bits for a trace, and \code{\#[derive(Trace)]}
registers every field of a unit under its field name; both read the
item's generic parameters so a generic unit can derive them.
\code{\#[pipeline(op = latency, ..)]} is the first of the lowering: on
an \code{async fn} of awaited operator calls it keeps the function and
emits a \code{NAME\_verilog()} beside it, each operator a stage of the
stated latency and each value delayed to meet the operator that uses
it. Its grammar is the one \code{mac} needs and no wider, and it
refuses the rest with its own message. \code{\#[lower]} on an
\code{impl Unit for Unit} does the same for a unit with state,
filling the header's \code{<In, Out>} in from \code{run}'s signature,
so the ports are named once: it reads \code{run} as one loop, or
several under
\code{join2}, each a process that waits for the rising or the falling
edge of its clock, reads registers and ports into names, predicates
drives with
\code{with!}, \code{case!} and \code{if}, chooses values
with
\code{select!},
and drives registers, words of memories and output ports, and writes
\code{Unit::lowered(name)}, \code{verilog} and \code{vhdl} beside
the impl. Ports come from \code{run}'s signature, a side of several
written as a tuple; registers, memories and widths come from the
struct through the \code{Trace} derive, which also derives
\code{Fields} and \code{Port}; and the configuration's constants are
evaluated when \code{lowered} runs, so a generic unit lowers once per
build. A \code{let} that computes is a wire named for it, a read of a
port or register an alias. It reads \code{case!} and \code{select!}
as chains on the value, with an enum variant, an integer or \code{\_}
as a pattern and a guard joined by \code{\&\&}; \code{slice},
\code{concat}, \code{sext} and \code{zext} by their turbofish, every
width it needs written, since it sees tokens and not types, and only
\code{concat}'s low operand left to Rust as \code{\_}; a
struct literal as the concatenation of its fields; an \code{if} chain
as arms under their conditions, each arm's statements lowered in
turn, a \code{send} in an arm hoisted with the arm's path as
\code{valid}, and a \code{set} of an output under it refused, since
an output is a wire; \code{self.m.at(a)}
on the left of a drive as a word of a memory and \code{self.m.read(a)}
as an index; and channels as a valid/ready handshake: \code{until} or
a channel's \code{wait} is the guard on every register drive after
it, a receive asserts \code{ready} under the guard, a \code{take}
binds \code{valid} and \code{data} to its two names and asserts
\code{ready} the same way, a send drives
\code{data} and asserts \code{valid} under it. What it does not read
it refuses with a message that names the construct.
A field of a port's value, \code{p.tag} where \code{p} is what a
channel port carries, is read as the slice of the port's data the
field occupies: the \code{Value} derive lays a struct's fields out
first field highest and says so in a \code{layout}, and the netlist
finds the bits when \code{lowered} runs, so a generic field's width
is no obstacle. \code{station!(Name, N)} writes a reservation
station of \code{N} inputs as text, a line struct, the station's
struct and its lowered \code{run}, every wire of the step named for
what it is, and \code{TXHDL\_MACRO\_DUMP} writes the text out for
reading; the attribute in its output expands as any does.
\code{interface!} takes an interface and any number of roles, checks
that every role names every member once and that no member is driven
twice, and emits a struct per role plus a constructor that splits each
member. \code{with!(self <= \{ .. \})} is the drives of one struct,
its name written once: an entry \code{field: value} is a drive,
\code{c ? field: value} a drive under a condition, \code{c ? \{ ..
\} else \{ .. \}} a group under one with the entries after
\code{else} under its failure, and a memory word \code{m.at(i)} a
field like any other; entries apply in order, the last drive of a
field winning, and each lowers as an \code{if} in the clocked block;
an entry that is not \code{field: value} is refused with the macro's
own message. \code{when!(c => self \{ .. \} else \{ .. \})} is one
group of it with the condition written first, the same entries and
the same \code{if}; an \code{if} whose condition is a constant of
the build is folded, its arm kept or dropped and the test gone. \code{case!} is a predicated drive per arm that reads
as a conditional: each statement inside an arm is \code{register <=
value}, the arrow pointing into the register, and a statement that
is not is refused with the macro's own message. Both are procedural
because a \code{macro\_rules!} \code{expr} fragment may be followed
only by \code{=>}, \code{,} or \code{;}, so \code{lhs <= rhs} cannot be
matched declaratively at all.

\lstinputlisting[caption={\code{lib/macros/lib.rs}.},label={lst:r-macros}]{lib/macros/lib.rs}
